% sections/interets.tex
% Centres d'Intérêt — montrer équilibre, esprit d'équipe, engagement dans la durée.
% Règle : chaque item inclut un quantificateur ou accomplissement concret.
% Éviter les hobbies passifs sans contexte (lecture, cinéma sans détail).

\section{Centres d'Intérêt}

\cvitem{Leadership \& Communication}{%
  Conférencier TEDx (Tec de Monterrey SLP, oct. 2022) devant +300 personnes ; discours publié sur le site officiel TEDx.
  Fondateur du groupe de mathématiques ROOT [sept. 2021 -- jan. 2023] :
  enseignement de la théorie des nombres à des lycéens avec distinctions étatiques aux Olympiades Mexicaines de Mathématiques (OMM) et au concours Gau55.%
}

\cvitem{Engagement communautaire}{%
  Mentor Technovation Girls [2025] : encadrement d'une équipe de 2 étudiantes en développement d'application.
  Bénévole Robotics WinterFest [2023, 2024] : co-organisation d'ateliers de robotique pratiques avec Robotistas.
  Bénévole Fundación Forge [mars 2026] : accompagnement de jeunes vers leur premier emploi.%
}

\cvitem{Musique}{%
  Second pianiste sur la production musicale \textit{Mamma Mia!} (campus Tec de Monterrey, nov. 2023) ;
  présentation Journée de la Guitare [oct. 2022, certificat à fournir].%
}

\cvitem{Compétition}{%
  Joueur d'échecs compétitif [confirmer tournois et années, inclure certificats] ;
  Membre de l'équipe robotique universitaire Tec de Monterrey (véhicules autonomes, vision par ordinateur).%
}

% TODO: Dès réception des certificats d'échecs, supprimer la note "[confirmer tournois et années]".
% TODO: Évaluer si la section dépasse la contrainte 1 page après compilation et réduire si besoin
%       en supprimant le cvitem "Compétition" ou en comprimant "Engagement communautaire" à 2 items.
